% ============================================================
% f2012_hubei_science_q19_fig2.tex
% 2012 湖北卷（理科）第 19 题 图 2（沿 AD 折起后的三棱锥 A-BCD）
% AD ⊥ 底面 BCD，∠BDC=90°；M 为 AC 中点，E 为 BC 中点；N 待定，原页未画出。
% 采用固定斜投影：D=(1.00,1.05)、B=(0.10,0.90)、C=(2.60,0.95)、A=(1.00,3.75)，
% 其中 E、M 由相应顶点直接取中点得到。
% 来源 PDF：普通高考/2012/2012湖北理.pdf（SHA-256 b84a88bc…07cc）
% 原页：第 2 页（A3 横排）第 1 栏下部，印面 738，第 19 题题干下方
% ============================================================
\begin{tikzpicture}[>=Stealth, line width=0.7pt, font=\small, x=1cm, y=1cm]
  \coordinate (A) at (1.00,3.75);
  \coordinate (B) at (0.10,0.90);
  \coordinate (C) at (2.60,0.95);
  \coordinate (D) at (1.00,1.05);
  \coordinate (E) at (1.35,0.925);
  \coordinate (M) at (1.80,2.35);

  % 可见棱
  \draw (A) -- (B);
  \draw (A) -- (C);
  \draw (B) -- (C);
  \draw (B) -- (M);

  % 被遮棱
  \draw[dashed] (A) -- (D);
  \draw[dashed] (D) -- (B);
  \draw[dashed] (D) -- (C);

  % D 处直角标记（DA 与 DC 之间）
  \draw (1.00,1.25) -- (1.20,1.2375) -- (1.20,1.0375);

  \fill (A) circle (1.2pt) node[above] {$A$};
  \fill (B) circle (1.2pt) node[below left] {$B$};
  \fill (C) circle (1.2pt) node[right] {$C$};
  \fill (D) circle (1.2pt) node[below left] {$D$};
  \fill (E) circle (1.2pt) node[below right] {$E$};
  \fill (M) circle (1.2pt) node[right] {$M$};
\end{tikzpicture}
